% !TeX spellcheck = en-US
% LTeX: language=en-US
% !TeX encoding = utf8
% !TeX TXS-program:compile = txs:///lualatex/[--shell-escape]
% !BIB program = biber
% -*- coding:utf-8 mod:LaTeX -*-

% The following package allows \\ at the title page
% For more information see https://github.com/latextemplates/scientific-thesis-cover/issues/4
\RequirePackage{kvoptions-patch}
\documentclass[
  % fontsize=11pt is the standard
  numbers=noenddot,
  english,  % English as main language; this parameter is passed to other packages (e.g., selnolig in the case of lualatex)
  a4paper,  % KOMAScript allows for both paper=a4 and (standard) a4paper - https://tex.stackexchange.com/a/61044/9075
  oneside,  % Changed from twoside to oneside to avoid jumping margins and page numbers for screen viewing
  bibliography=totoc,
  % idxtotoc,   % Index ins Inhaltsverzeichnis
  % liststotoc, % List of * ins Inhaltsverzeichnis, mit liststotocnumbered werden die Abbildungsverzeichnisse nummeriert
  headsepline,
  cleardoublepage=empty,
  parskip=half,
  %               draft    % um zu sehen, wo noch nachgebessert werden muss - wichtig, da Bindungskorrektur mit drin
  draft=false
]{scrbook}
\usepackage{scrlayer-scrpage}

\usepackage{iftex}

\usepackage{ifplatform}

% backticks (`) are rendered as such in verbatim environments.
% See following links for details:
%   - https://tex.stackexchange.com/a/341057/9075
%   - https://tex.stackexchange.com/a/47451/9075
%   - https://tex.stackexchange.com/a/166791/9075
\usepackage{upquote}

% Set English as language and allow to write hyphenated"=words
%
% Even though `american`, `english` and `USenglish` are synonyms for babel package (according to https://tex.stackexchange.com/questions/12775/babel-english-american-usenglish), the llncs document class is prepared to avoid the overriding of certain names (such as "Abstract." -> "Abstract" or "Fig." -> "Figure") when using `english`, but not when using the other 2.
% english has to go last to set it as default language
\usepackage[ngerman,main=english]{babel}
%
% Hint by http://tex.stackexchange.com/a/321066/9075 -> enable "= as dashes
\addto\extrasenglish{\languageshorthands{ngerman}\useshorthands{"}}

% Links behave as they should. Enables "\url{...}" for URL typesettings.
% Allow URL breaks also at a hyphen, even though it might be confusing: Is the "-" part of the address or just a hyphen?
% See https://tex.stackexchange.com/a/3034/9075.
\usepackage[hyphens]{url}

% When activated, use text font as url font, not the monospaced one.
% For all options see https://tex.stackexchange.com/a/261435/9075.
% \urlstyle{same}

% Improve wrapping of URLs - hint by http://tex.stackexchange.com/a/10419/9075
\makeatletter
\g@addto@macro{\UrlBreaks}{\UrlOrds}
\makeatother

% nicer // - solution by http://tex.stackexchange.com/a/98470/9075
% DO NOT ACTIVATE -> prevents line breaks
%\makeatletter
%\def\Url@twoslashes{\mathchar`\/\@ifnextchar/{\kern-.2em}{}}
%\g@addto@macro\UrlSpecials{\do\/{\Url@twoslashes}}
%\makeatother

%math stuff
\usepackage[
  centertags,    % (default) center tags vertically
  % tbtags,        % 'Top-or-bottom tags': For a split equation, place equation numbers level with the last (resp. first) line, if numbers are on the right (resp. left).
  sumlimits,    % (default) Place the subscripts and superscripts of summation symbols above and below
  % nosumlimits,   % Always place the subscripts and superscripts of summation-type symbols to the side, even in displayed equations.
  intlimits,     % Like sumlimits, but for integral symbols.
  % nointlimits,   % (default) Opposite of intlimits.
  namelimits,    % (default) Like sumlimits, but for certain 'operator names' such as det, inf, lim, max, min, that traditionally have subscripts placed underneath when they occur in a displayed equation.
  % nonamelimits,  % Opposite of namelimits.
  % leqno,         % Place equation numbers on the left.
  % reqno,         % Place equation numbers on the right.
  fleqn,         % Position equations at a fixed indent from the left margin rather than centered in the text column.
]{amsmath}
\SetMathAlphabet{\mathcal}{normal}{OMS}{amsa}{m}{n} %% AMS font for mathcal

%%% Doc: http://mirror.ctan.org/tex-archive/macros/latex/contrib/mh/doc/mathtools.pdf
% Erweitert amsmath und behebt einige Bugs
\usepackage[fixamsmath,disallowspaces]{mathtools}

%%% Doc: http://www.ctan.org/info?id=fixmath
% LaTeX's default style of typesetting mathematics does not comply
% with the International Standards ISO31-0:1992 to ISO31-13:1992
% which indicate that uppercase Greek letters always be typeset
% upright, as opposed to italic (even though they usually
% represent variables) and allow for typesetting of variables in a
% boldface italic style (even though the required fonts are
% available). This package ensures that uppercase Greek be typeset
% in italic style, that upright $\Delta$ and $\Omega$ symbols are
% available through the commands \upDelta and \upOmega; and
% provides a new math alphabet \mathbold for boldface
% italic letters, including Greek.
\usepackage{fixmath}

%for theorems, replacement for amsthm
\usepackage[amsmath,hyperref]{ntheorem}
\theorempreskipamount 2ex plus1ex minus0.5ex
\theorempostskipamount 2ex plus1ex minus0.5ex
\theoremstyle{break}
\newtheorem{definition}{Definition}[chapter]

%%% Doc: http://mirror.ctan.org/tex-archive/macros/latex/contrib/onlyamsmath/onlyamsmath.dvi
% Warnt bei Benutzung von Befehlen die mit amsmath inkompatibel sind.

% Braucht man evtl. nicht.
% \usepackage[
% 	all,
% 	warning
% ]{onlyamsmath}

%% !!! If you change the font, be sure that words such as "workflow" can
%% !!! still be copied from the PDF. If this is not the case, you have
%% !!! to use glyphtounicode. See comment at cmap package.
%%
%% Background: "workflow" contains "fl" which is a ligature, which in turn
%%             is rendered as one character in the PDF and needs to be split
%%             whily copying.

\ifluatex
  \usepackage[no-math]{fontspec}
  \usepackage{unicode-math}

  % See https://tug.org/FontCatalogue/texgyretermes/ for more information
  \setmainfont{texgyretermes}[
    Extension = .otf,
    UprightFont = *-regular,
    BoldFont = *-bold,
    ItalicFont = *-italic,
    BoldItalicFont = *-bolditalic,
    Ligatures=TeX
  ]
  % Use file-based loading for TeX Gyre Heros to avoid OS font dependency
  \setsansfont{texgyreheros}[
    Scale=.9,
    Extension = .otf,
    UprightFont = *-regular,
    BoldFont = *-bold,
    ItalicFont = *-italic,
    BoldItalicFont = *-bolditalic
  ]
  % Use the default Latin Modern Mono instead of Inconsolata to avoid missing font errors
  % \setmonofont[StylisticSet={1,3},Scale=.9]{Inconsolata}

  % Enable proper ligatures
  % For more information see https://ctan.org/pkg/selnolig
  % language "english" or "ngerman" is passed to selnolig by the document class
  \usepackage{selnolig}

\else
  \RequirePackage{newtxtext}
  \RequirePackage{newtxmath}
  \RequirePackage[zerostyle=b,scaled=.9]{newtxtt}

  % Has to be loaded AFTER any font packages. See https://tex.stackexchange.com/a/2869/9075.
  \usepackage[T1]{fontenc}
\fi

% DE: Noch mehr Symbole
%\usepackage{stmaryrd} %fuer \ovee, \owedge, \otimes
%\usepackage{marvosym} %fuer \Writinghand %patched to not redefine \Rightarrow
%\usepackage{mathrsfs} %mittels \mathscr{} schoenen geschwungenen Buchstaben erzeugen
%\usepackage{calrsfs} %\mathcal{} ein bisserl dickeren buchstaben erzeugen - sieht net so gut aus.

% EN: Fallback font - if the subsequent font packages do not define a font (e.g., monospaced)
%     This is the modern package for "Computer Modern".
%     In case this gets activated, one has to switch from cmap package to glyphtounicode (in the case of pdflatex)
% DE: Fallback-Schriftart
%\usepackage[%
%    rm={oldstyle=false,proportional=true},%
%    sf={oldstyle=false,proportional=true},%
%    tt={oldstyle=false,proportional=true,variable=true},%
%    qt=false%
%]{cfr-lm}

% EN: Headings are typeset in Helvetica (which is similar to Arial)
% DE: Schriftart fuer die Ueberschriften - ueberschreibt lmodern
%\usepackage[scaled=.95]{helvet}

% DE: Für Schreibschrift würde tun, muss aber nicht
%\usepackage{mathrsfs} %  \mathscr{ABC}

% EN: Font for the main text
% DE: Schriftart fuer den Fliesstext - ueberschreibt lmodern
%     Linux Libertine, siehe http://www.linuxlibertine.org/
%     Packageparamter [osf] = Minuskel-Ziffern
%     rm = libertine im Brottext, Linux Biolinum NICHT als serifenlose Schrift, sondern helvet (von oben) beibehalten
%\usepackage[rm]{libertine}

% EN: Alternative Font: Palantino. It is recommeded by Prof. Ludewig for German texts
% DE: Alternative Schriftart: Palantino, Packageparamter [osf] = Minuskel-Ziffern
%     Bitte nur in deutschen Texten
%\usepackage{mathpazo} %ftp://ftp.dante.de/tex-archive/fonts/mathpazo/ - Tipp aus DE-TEX-FAQ 8.2.1
% EN: The euro sign
% DE: Das Euro Zeichen
%     Fuer Palatino (mathpazo.sty): richtiges Euro-Zeichen
%     Alternative: \usepackage{eurosym}
% \newcommand{\EUR}{\ppleuro}

% DE: Schriftart fuer Programmcode - ueberschreibt lmodern
%     Falls auskommentiert, wird die Standardschriftart lmodern genommen
%     Fuer schreibmaschinenartige Schluesselwoerter in den Listings - geht bei alten Installationen nicht, da einige Fontshapes (<>=) fehlen
%\usepackage[scaled=.92]{luximono}
%\usepackage{courier}
% DE: BeraMono als Typewriter-Schrift, Tipp von http://tex.stackexchange.com/a/71346/9075
%\usepackage[scaled=0.83]{beramono}

\usepackage{setspace}
% Alternative package: https://ctan.org/pkg/leading

% Symbole Check und Cross
\usepackage{pifont}
\newcommand{\dingcheck}{\ding{51}}
\newcommand{\dingcross}{\ding{55}}
%for scaling see http://tex.stackexchange.com/a/130236/9075

% DE: Noch mehr Symbole
%\usepackage{stmaryrd} %fuer \ovee, \owedge, \otimes
%\usepackage{marvosym} %fuer \Writinghand %patched to not redefine \Rightarrow
%\usepackage{mathrsfs} %mittels \mathscr{} schoenen geschwungenen Buchstaben erzeugen
%\usepackage{calrsfs} %\mathcal{} ein bisserl dickeren buchstaben erzeugen - sieht net so gut aus.

\automark[section]{chapter}
\setkomafont{pageheadfoot}{\normalfont\sffamily}
\setkomafont{pagenumber}{\normalfont\sffamily}

% Force paragraph headings to have a normal space instead of a huge space
\RedeclareSectionCommand[afterskip=-0.5em]{paragraph}

\ihead[]{}
\chead[]{}
\ohead[]{\headmark}
\cfoot[]{}
\ofoot[\usekomafont{pagenumber}\thepage]{\usekomafont{pagenumber}\thepage}
\ifoot[]{}

% Character protrusion and font expansion. See http://www.ctan.org/tex-archive/macros/latex/contrib/microtype/

\usepackage[
  babel=true, % Enable language-specific kerning. Take language-settings from the languge of the current document (see Section 6 of microtype.pdf)
  expansion=alltext,
  protrusion=alltext-nott, % Ensure that at listings, there is no change at the margin of the listing
  % In the standard configuration, this template is always in the final mode, so this option only makes a difference if "pros" use the draft mode
  final % Always enable microtype, even if in draft mode. This helps finding bad boxes quickly.
]{microtype}

% \texttt{test -- test} keeps the "--" as "--" (and does not convert it to an en dash)
\DisableLigatures{encoding = T1, family = tt* }

%\DeclareMicrotypeSet*[tracking]{my}{ font = */*/*/sc/* }%
%\SetTracking{ encoding = *, shape = sc }{ 45 }
% Source: http://homepage.ruhr-uni-bochum.de/Georg.Verweyen/pakete.html
% Deactiviated, because does not look good

\usepackage{graphicx}

\usepackage{tikz}
\usetikzlibrary{positioning, arrows.meta, shapes.geometric, backgrounds, shapes.symbols}

% Base folder, so there is no need to repeat this over and over again.
\graphicspath{ {figures/} }

%%% Doc: http://mirror.ctan.org/tex-archive/macros/latex/contrib/pdfpages/pdfpages.pdf
\usepackage{pdfpages} % Include pages from external PDF documents in LaTeX documents

% Diagonal lines in a table - http://tex.stackexchange.com/questions/17745/diagonal-lines-in-table-cell
% Slashbox is not available in texlive (due to licensing) and also gives bad results. Thus, we use diagbox
\usepackage{diagbox}

\ifluatex
  \usepackage{spelling}
  \spellingoutput{off}
\fi

\usepackage[dvipsnames, table]{xcolor}
% Code Listings
\usepackage{listings}

\definecolor{eclipseStrings}{RGB}{42,0.0,255}
\definecolor{eclipseKeywords}{RGB}{127,0,85}
\colorlet{numb}{magenta!60!black}

% JSON definition
% Source: https://tex.stackexchange.com/a/433961/9075

\lstdefinelanguage{json}{
  basicstyle=\normalfont\ttfamily,
  commentstyle=\color{eclipseStrings}, % style of comment
  stringstyle=\color{eclipseKeywords}, % style of strings
  numbers=left,
  numberstyle=\scriptsize,
  stepnumber=1,
  numbersep=8pt,
  showstringspaces=false,
  breaklines=true,
  frame=lines,
  % backgroundcolor=\color{gray}, %only if you like
  string=[s]{"}{"},
  comment=[l]{:\ "},
  morecomment=[l]{:"},
  literate=
    *{0}{{{\color{numb}0}}}{1}
    {1}{{{\color{numb}1}}}{1}
    {2}{{{\color{numb}2}}}{1}
    {3}{{{\color{numb}3}}}{1}
    {4}{{{\color{numb}4}}}{1}
    {5}{{{\color{numb}5}}}{1}
    {6}{{{\color{numb}6}}}{1}
    {7}{{{\color{numb}7}}}{1}
    {8}{{{\color{numb}8}}}{1}
    {9}{{{\color{numb}9}}}{1}
}

% PDDL language definition for planning domain listings
\lstdefinelanguage{PDDL}{
  morekeywords={define, domain, problem, requirements, strips, typing, equality,
    predicates, action, parameters, precondition, effect, objects, init, goal,
    and, not, or, types, constants, functions, either, forall, exists, when, imply},
  sensitive=false,
  morecomment=[l]{;},
}

\lstset{
  % everything between (* *) is a latex command
  escapeinside={(*}{*)},
  %
  language=json,
  %
  showstringspaces=false,
  %
  basicstyle=\footnotesize\ttfamily,
  %
  commentstyle=\slshape,
  %
  % default: \rmfamily
  stringstyle=\ttfamily,
  %
  breaklines=true,            % Zeilen werden umbrochen
  %
  breakatwhitespace=true,
  %
  % alternative: fixed
  columns=flexible,
  %
  tabsize=2,                  % Groesse von Tabs
  %
  numbers=left,
  %
  numberstyle=\tiny,
  %
  basewidth=.5em,
  %
  xleftmargin=.5cm,
  %
  % aboveskip=0mm,
  %
  % belowskip=0mm,
  %
  captionpos=b
}

\ifpdftex
  % Enable Umlauts when using \lstinputputlisting.
  % See https://stackoverflow.com/a/29260603/873282 and https://tex.stackexchange.com/a/24532/9075 for details.
  % listingsutf8 did not work in June 2020.
  \lstset{extendedchars=true, literate=
    {á}{{\'a}}1 {é}{{\'e}}1 {í}{{\'i}}1 {ó}{{\'o}}1 {ú}{{\'u}}1
  {Á}{{\'A}}1 {É}{{\'E}}1 {Í}{{\'I}}1 {Ó}{{\'O}}1 {Ú}{{\'U}}1
  {à}{{\`a}}1 {è}{{\`e}}1 {ì}{{\`i}}1 {ò}{{\`o}}1 {ù}{{\`u}}1
  {À}{{\`A}}1 {È}{{\'E}}1 {Ì}{{\`I}}1 {Ò}{{\`O}}1 {Ù}{{\`U}}1
  {ä}{{\"a}}1 {ë}{{\"e}}1 {ï}{{\"i}}1 {ö}{{\"o}}1 {ü}{{\"u}}1
  {Ä}{{\"A}}1 {Ë}{{\"E}}1 {Ï}{{\"I}}1 {Ö}{{\"O}}1 {Ü}{{\"U}}1
  {â}{{\^a}}1 {ê}{{\^e}}1 {î}{{\^i}}1 {ô}{{\^o}}1 {û}{{\^u}}1
  {Â}{{\^A}}1 {Ê}{{\^E}}1 {Î}{{\^I}}1 {Ô}{{\^O}}1 {Û}{{\^U}}1
  {Ã}{{\~A}}1 {ã}{{\~a}}1 {Õ}{{\~O}}1 {õ}{{\~o}}1
  {œ}{{\oe}}1 {Œ}{{\OE}}1 {æ}{{\ae}}1 {Æ}{{\AE}}1 {ß}{{\ss}}1
  {ű}{{\H{u}}}1 {Ű}{{\H{U}}}1 {ő}{{\H{o}}}1 {Ő}{{\H{O}}}1
  {ç}{{\c c}}1 {Ç}{{\c C}}1 {ø}{{\o}}1 {å}{{\r a}}1 {Å}{{\r A}}1
  }
\fi

\lstloadlanguages{% Check dokumentation for further languages...
  %[Visual]Basic
  %Pascal
  %C
  %C++
  %XML
  %HTML
}

% For easy quotations: \enquote{text}
% This package is very smart when nesting is applied, otherwise textcmds (see below) provides a shorter command
\usepackage[autostyle=true]{csquotes}

% Enable using "`quote"' - see https://tex.stackexchange.com/a/150954/9075
\defineshorthand{"`}{\openautoquote}
\defineshorthand{"'}{\closeautoquote}

% Nicer tables (\toprule, \midrule, \bottomrule)
\usepackage{booktabs}

% Extended enumerate, such as \begin{compactenum}
\usepackage{paralist}
\usepackage[
  backend       = biber, %biber does not work with 64x versions alternative: bibtex8; minalphanames only works with biber backend
  sortcites     = true,
  bibstyle      = alphabetic,
  citestyle     = alphabetic,
  giveninits    = true,
  useprefix     = false, %"von, van, etc." will be printed, too. See below.
  minnames      = 1,
  minalphanames = 3,
  maxalphanames = 4,
  maxbibnames   = 99,
  maxcitenames  = 2,
  natbib        = true,
  eprint        = true,
  url           = true,
  doi           = true, %source: http://tex.stackexchange.com/a/23118/9075
  isbn          = true, %source: http://tex.stackexchange.com/a/23118/9075
  backref       = true]{biblatex}

% enable more breaks at URLs. See https://tex.stackexchange.com/a/134281.
\setcounter{biburllcpenalty}{7000}
\setcounter{biburlucpenalty}{8000}

\bibliography{bibliography}
%\addbibresource[datatype=bibtex]{\bibliography{bibliography}}

% Do not put "vd" in the label, but put it at "\citeauthor"
% Source: http://tex.stackexchange.com/a/30277/9075
\makeatletter
\AtBeginDocument{\toggletrue{blx@useprefix}}
\AtBeginBibliography{\togglefalse{blx@useprefix}}
\makeatother

% Thin spaces between initials
% http://tex.stackexchange.com/a/11083/9075
\renewrobustcmd*{\bibinitdelim}{\,}

% Keep first and last name together in the bibliography
% http://tex.stackexchange.com/a/196192/9075
\renewcommand*\bibnamedelimc{\addnbspace}
\renewcommand*\bibnamedelimd{\addnbspace}

% Replace last "and" by comma in bibliography
% See http://tex.stackexchange.com/a/41532/9075
\AtBeginBibliography{%
  \renewcommand*{\finalnamedelim}{\addcomma\space}%
}

% enable hyperlinked author names when using \citeauthor
% source: http://tex.stackexchange.com/a/75916/9075
\DeclareCiteCommand{\citeauthor}
  {\boolfalse{citetracker}%
   \boolfalse{pagetracker}%
   \usebibmacro{prenote}}
  {\ifciteindex
     {\indexnames{labelname}}
     {}%
   \printtext[bibhyperref]{\printnames{labelname}}}
  {\multicitedelim}
  {\usebibmacro{postnote}}

% Farbige Tabellen
% ----------------
% Das Paket colortbl wird inzwischen automatisch durch xcolor geladen
%
% Erweiterte Funktionen innerhalb von Tabellen
% --------------------------------------------
%%% Doc: http://mirror.ctan.org/tex-archive/macros/latex/contrib/multirow/multirow.sty
\usepackage{multirow} % Mehrfachspalten
%
%%% Doc: Documentation inside dtx Package
\usepackage{dcolumn}  % Ausrichtung an Komma oder Punkt

%%% Doc: http://mirror.ctan.org/tex-archive/macros/latex/contrib/supertabular/supertabular.pdf
%\usepackage{supertabular}

%%% Fussnoten/Endnoten ===================================================

%%% Doc: http://mirror.ctan.org/tex-archive/macros/latex/contrib/footmisc/footmisc.pdf
%
\usepackage[
  bottom,      % Footnotes appear always on bottom. This is necessary specially when floats are used
  stable,      % Make footnotes stable in section titles
  % perpage,     % Reset on each page
  % para,        % Place footnotes side by side of in one paragraph.
  % side,        % Place footnotes in the margin
  ragged,      % Use RaggedRight
  % norule,      % Suppress rule above footnotes
  multiple,    % Rearrange multiple footnotes intelligent in the text.
  % symbol,      % Use symbols instead of numbers
]{footmisc}

\counterwithout{footnote}{chapter} % Continuous numbering of footnotes across chapters

\interfootnotelinepenalty=10000 % Verhindert das Fortsetzen von Fussnoten auf der gegenüberligenden Seite

% EN: Put footnotes below floats
% DE: Fußnoten unter Gleitumgebungen ("floats") platzieren
% Source: https://tex.stackexchange.com/a/32993/9075
\usepackage{stfloats}
\fnbelowfloat

% EN: Extended support for footnotes
% DE: Fußnoten
%
%\usepackage{dblfnote}  %Zweispaltige Fußnoten
%
% Keine hochgestellten Ziffern in der Fußnote (KOMA-Script-spezifisch):
%\deffootnote[1.5em]{0pt}{1em}{\makebox[1.5em][l]{\bfseries\thefootnotemark}}
%
% Abstand zwischen Fußnoten vergrößern:
%\setlength{\footnotesep}{.85\baselineskip}
%
% EN: Following command disables the separting line of the footnote
% DE: Folgendes Kommando deaktiviert die Trennlinie zur Fußnote
%\renewcommand{\footnoterule}{}
%
%\addtolength{\skip\footins}{\baselineskip} % Abstand Text <-> Fußnote

% DE: Fußnoten immer ganz unten auf einer \raggedbottom-Seite
% DE: fnpos kommt aus dem yafoot package
%\usepackage{fnpos}
%\makeFNbelow
%\makeFNbottom

% TODO (and comment) configuration
%
% - \todo (from todo, easy-todo, todonotes) / \TODO (from fixmetodonotes) - for "normal" TODOs
% - \todofix - "important" TODOs
%
% - \textcomment - highlights text and has a hover comment
% - \sidecomment - just puts a comment to the side. Note: \comment MUST NOT be used as command name, it is already defined by much packages (mathdesign, mindflow, verbatim, and others)
%
% - \missingfigure
%
% - \textmarker
% - \modified
% - \change      - adresses a review comment

% Enable nice comments
\usepackage{pdfcomment}

%\newcommand{\textcomment}[2]{\colorbox{yellow!60}{#1}\pdfcomment[color={0.234 0.867 0.211},hoffset=-6pt,voffset=10pt,opacity=0.5]{#2}}
\newcommand{\textcomment}[2]{#1}

% Small PDF comment
% 1. Parameter: Comment
%\newcommand{\sidecomment}[1]{\pdfcomment[color={0.045 0.278 0.643},voffset=4pt,icon=Note]{#1}}
% Disabled variant - for the final PDF
\newcommand{\sidecomment}[1]{}

\newcommand{\todo}[1]{TODO!\sidecomment{#1}}

% Änderungen
%
% 1. Parameter: Review-Kommentar
% 2. Parameter: Neuer Text
%\newcommand{\change}[2]{{\color{red}#2}\pdfcomment[color={0.234 0.867 0.211},voffset=8pt,opacity=0.5]{#1}}
% Disabled variant - for the final PDF
\newcommand{\change}[2]{#2}

% Define default commands
\makeatletter
\@ifundefined{missingfigure}{\newcommand{\missingfigure}{... missing figure ...}}{}
\@ifundefined{textcomment}{\newcommand{\textcomment}[2]{#1 \todo{#2}}}{}
\@ifundefined{sidecomment}{\newcommand{\sidecomment}[1]{\marginpar{#1}}}{}
\@ifundefined{todo}{\newcommand{\todo}[1]{\sidecomment{#1}}}{}
\@ifundefined{TODO}{\newcommand{\TODO}[1]{\todo{#1}}}{}
\@ifundefined{todofix}{\newcommand{\todofix}[1]{\todo{#1}}}{}
\@ifundefined{change}{\newcommand{\change}[2]{#1 $\rightarrow$ #2}}{}
\makeatother

% Textmarker (Textfarbe rot)
\newcommand{\textmarker}[1]{{\color{red} #1}\xspace}

% Modified (Text blau)
\newcommand{\modified}[1]{{\color{blue!60!black} #1}\xspace}

\usepackage[group-minimum-digits=4,per-mode=fraction]{siunitx}

% See http://tex.stackexchange.com/a/83051/9075
% Normally, doesn't work with hyperref, but cleveref fixes that
\usepackage{varioref}

% Enable that parameters of \cref{}, \ref{}, \cite{}, ... are linked so that a reader can click on the number an jump to the target in the document
\usepackage{hyperref}

% Enable hyperref without colors and without bookmarks
\hypersetup{
  hidelinks,
  colorlinks=true,       % Links erhalten Farben statt Kaeten
  raiselinks=true,       % calculate real height of the link
  allcolors=black,
  pdfstartview=Fit,
  breaklinks=true,       % Links ueberstehen Zeilenumbruch
  hypertexnames=false,   % Fix jumping to algorithm line - http://tex.stackexchange.com/a/156404/9075
}

% Enable correct jumping to figures when referencing
\usepackage[all]{hypcap}


%%%
% Ermoeglicht es, Abbildungen um 90 Grad zu drehen
% Alternatives Paket: rotating Allerdings wird hier nur das Bild gedreht, während bei lscape auch die PDF-Seite gedreht wird.
%Das Paket lscape dreht die Seite auch nicht
\usepackage{pdflscape}

\usepackage[caption=false,font=footnotesize]{subfig}

% Alternative for making subfigures:
% Part of the caption package. See http://www.ctan.org/pkg/caption
% Ersetzt die Pakete subfigure und subfig - siehe https://tex.stackexchange.com/a/13778/9075
%
% (subfigure is outdated. subfig is maintained, but subcaption is better)
% See: http://tex.stackexchange.com/questions/13625/subcaption-vs-subfig-best-package-for-referencing-a-subfigure
%\usepackage[hypcap=true]{subcaption}

\usepackage{mindflow}

% https://ctan.org/pkg/algorithms
% Consists of two environments: algorithm and algorithmic
% Although oudated, it defines the "algorithm" float enviornment
% TODO: Define floating environment "algorithm" in other ways
\usepackage[chapter]{algorithm}

% https://ctan.org/pkg/algpseudocodex
% Successor of algorithmicx; more modern than https://ctan.org/pkg/algorithms
\usepackage{algpseudocodex}

\floatname{algorithm}{Algorithmus}
\renewcommand{\listalgorithmname}{Algorithmenverzeichnis}

\newcommand{\commentchar}{\ensuremath{/\mkern-4mu/}}
\algrenewcommand{\algorithmiccomment}[1]{\hfill $\commentchar$ #1}

% Extensions for references inside the document (\cref{fig:sample}, ...)
% Enable usage \cref{...} and \Cref{...} instead of \ref: Type of reference included in the link
% That means, "Figure 5" is a full link instead of just "5".
\usepackage[capitalise,nameinlink,noabbrev]{cleveref}

\crefname{listing}{Listing}{Listings}
\Crefname{listing}{Listing}{Listings}
\crefname{lstlisting}{Listing}{Listings}
\Crefname{lstlisting}{Listing}{Listings}

\usepackage{lipsum}

% For demonstration purposes only
% These packages can be removed when all examples have been deleted
\usepackage[math]{blindtext}
\usepackage{mwe}
\usepackage[realmainfile]{currfile}
\usepackage{tcolorbox}
\tcbuselibrary{listings}

%introduce \powerset - hint by http://matheplanet.com/matheplanet/nuke/html/viewtopic.php?topic=136492&post_id=997377
\DeclareFontFamily{U}{MnSymbolC}{}
\DeclareSymbolFont{MnSyC}{U}{MnSymbolC}{m}{n}
\DeclareFontShape{U}{MnSymbolC}{m}{n}{
  <-6>    MnSymbolC5
  <6-7>   MnSymbolC6
  <7-8>   MnSymbolC7
  <8-9>   MnSymbolC8
  <9-10>  MnSymbolC9
  <10-12> MnSymbolC10
  <12->   MnSymbolC12%
}{}
\DeclareMathSymbol{\powerset}{\mathord}{MnSyC}{180}

\usepackage[
  translate=babel,
  abbreviations,         % create "abbreviations" glossary
  nomain,                % don't create "main" glossary
  stylemods=longbooktabs % do the adjustments for the longbooktabs styles
]{glossaries-extra}
\setglossarystyle{long3col-booktabs}

% Hint by https://tex.stackexchange.com/a/463188/9075
% \usepackage{glossary-longextra}

% Following is required if the abbreviation list should be sorted automatically (\printglossary / \printglossaries)
% Not required, if we printed the entries in-order (using \printunsrtglossaries)
% Required to have the German chapter name % Source: https://tex.stackexchange.com/a/426392/9075
\makeglossaries

\newabbreviation{ai}{AI}{Artificial Intelligence}
\newabbreviation{llm}{LLM}{Large Language Model}
\newabbreviation{api}{API}{Application Programming Interface}
\newabbreviation{pddl}{PDDL}{Planning Domain Definition Language}
\newabbreviation{ipc}{IPC}{International Planning Competition}
\newabbreviation{sat}{SAT}{Boolean Satisfiability}
\newabbreviation{bfws}{BFWS}{Best-First Width Search}
\newabbreviation{cot}{CoT}{Chain-of-Thought}
\newabbreviation{scc}{SCC}{Strongly Connected Component}
\newabbreviation{par10}{PAR10}{Penalised Average Runtime (factor~10)}
\newabbreviation{htn}{HTN}{Hierarchical Task Network}
\newabbreviation{strips}{STRIPS}{Stanford Research Institute Problem Solver}
\newabbreviation{adl}{ADL}{Action Description Language}


% Allows for defining commands that don't eat spaces.
\usepackage{xspace}
% Adds compatibility to \xspace und \enquote
\makeatletter
\xspaceaddexceptions{\grqq \grq \csq@qclose@i \} }
\makeatother

\newcommand{\eg}{e.g.,\ }
\newcommand{\ie}{i.e.,\ }

% Enable hyphenation at other places as the dash.
% Example: applicaiton\hydash specific
\makeatletter
\newcommand{\hydash}{\penalty\@M-\hskip\z@skip}
% Definition of "= taken from http://mirror.ctan.org/macros/latex/contrib/babel-contrib/german/ngermanb.dtx
\makeatother

% Add manual adapted hyphenation of English words
% See https://ctan.org/pkg/hyphenex and https://tex.stackexchange.com/a/22892/9075 for details
\input{ushyphex}

% correct bad hyphenation here
\hyphenation{
  op-tical net-works semi-conduc-tor
  % May not be hypphenated
  AROMA TOSCA BPMN OASIS OMG DMTF IT DevOps
}

\input{commands}

% Package URL: https://ctan.org/pkg/scientific-thesis-cover
\usepackage[
  title={Architecture-Aware Domain Model Configuration: Leveraging LLMs and Feedback Loops for AI Planner-Specific Optimization},
  author={Daniel Ehab Rasmy Bendas},
  type=bachelor,
  institute=iaas,
  course={Computer Science and Engineering},
  examiner={Dr.\ Ilche Georgievski},
  supervisor={Dr.\ Ilche Georgievski},
  startdate={March 1, 2026},
  enddate={July 1, 2026},
  language=english
]{scientific-thesis-cover}


\ifpdftex
  % Enable copy and paste of text from the PDF
  % Only required for pdflatex. It "just works" in the case of lualatex.
  % Alternative: cmap or mmap package
  % mmap enables mathematical symbols, but does not work with the newtx font set
  % See: https://tex.stackexchange.com/a/64457/9075
  % Other solutions outlined at http://goemonx.blogspot.de/2012/01/pdflatex-ligaturen-und-copynpaste.html and http://tex.stackexchange.com/questions/4397/make-ligatures-in-linux-libertine-copyable-and-searchable
  % Trouble shooting outlined at https://tex.stackexchange.com/a/100618/9075
  %
  % According to https://tex.stackexchange.com/q/451235/9075 this is the way to go
  \input{glyphtounicode}
  \pdfgentounicode=1
\fi
% DM: line-breaking-description env vom daniel w.

% credit goes to daniel w. :-)
%% --- Descriptions with line breaks in labels ---------------------------------
\usepackage{calc}

\newcommand*\Descriptionlabel[1]{%
  \raisebox{0pt}[1ex][0pt]{
    \makebox[\labelwidth][1]{
      \parbox[t]{\labelwidth}{
        \hspace{0pt}\textbf{#1:}}}}
}

\newcommand*\Descriptionlabelx[1]{%
  \parbox[t]{\textwidth}{
    \textbf{#1}\\\mbox{}}
}

\newenvironment{Description}{
  \begin{list}{}{
      \let\makelabel\Descriptionlabelx
      \setlength\labelwidth{1em}
      \setlength\leftmargin{\labelwidth+\labelsep}
    }
    }
    {
  \end{list}
}

% globally change line spacing of lists
% paralist has suspended development since 10 years.
% enumitem has been updated 2011-09-28
\usepackage[inline]{enumitem}
\setlist{partopsep=0pt,itemsep=1pt}

%------------------------------------------------------------------------
% fquote Fancy Quotation environment
% supports empty/optional author

% Use \sloppy to make right-margin easier?
% Set picture units to be relative to font size (em)?
% Use begingroup to rest units afterwards?

\usepackage{xifthen}% provides \isempty test
\definecolor{quotemark}{gray}{0.7}

%fquote environment with author as optional parameter
%usage: \begin{fquote}quote\end{fquote} or \begin{fquote}[Author]quote\end{fquote}
\newenvironment{fquote}[1][]{%
  \newcommand{\fqauthor}{\relax}
  \ifthenelse{\isempty{#1}}
  {}% do nothing
  {\renewcommand{\fqauthor}{\hfill\textsc{--- #1}}}
  \vspace{1em}
  \begin{list}{}{%
      \setlength{\leftmargin}{0.2\textwidth}
      \setlength{\rightmargin}{0.2\textwidth}
    }
    \item[]%
          \begin{picture}(0,0)(0,0)
            \put(-15,-5){\makebox(0,0){%
                \scalebox{4.5}{\textcolor{quotemark}{\bfseries``}}}%
            }
          \end{picture}\em\ignorespaces%
          }{%
          \newline%
          \makebox[0pt][l]{\hspace{0.6\textwidth}%
            \begin{picture}(0,0)(0,0)
              \put(15,10){\makebox(0,0){%
                  \scalebox{4.5}{\textcolor{quotemark}{\rmfamily\bfseries''}}}%
              }
            \end{picture}}%
          \fqauthor
  \end{list}
}

%German fquote
%  1 parameter for the author's name, may be empty ("{}")
%  guaranteed German quotes (works with lualatex and babel package)
%  usage: \begin{gfquote}{Author}quote\end{gfquote}
\newenvironment{gfquote}[1]{%
  \newcommand{\fqauthor}{\relax}
  \ifthenelse{\isempty{#1}}
  {}% do nothing
  {\renewcommand{\fqauthor}{\hfill\textsc{\textemdash #1}}}
  \vspace{1em}
  \begin{list}{}{%
      \setlength{\leftmargin}{0.2\textwidth}
      \setlength{\rightmargin}{0.2\textwidth}
    }
    \item[]%
          \begin{picture}(0,0)(0,0)
            \put(-15,-5){\makebox(0,0){%
                \scalebox{4.5}{\textcolor{quotemark}{\bfseries \glqq}}}%
            }
          \end{picture}\em\ignorespaces%
          }{%
          \newline%
          \makebox[0pt][l]{\hspace{0.6\textwidth}%
            \begin{picture}(0,0)(0,0)
              \put(15,10){\makebox(0,0){%
                  \scalebox{4.5}{\textcolor{quotemark}{\rmfamily\bfseries \grqq}}}%
              }
            \end{picture}}%
          \fqauthor
  \end{list}
}

% fix incompatibilities between KOMA and other packages, mainly float.
% should be loaded at the very end - see http://tex.stackexchange.com/a/156256/9075
\usepackage{scrhack}


\begin{document}
\raggedbottom
\pagenumbering{arabic}
\Titelblatt

\pagestyle{plain.scrheadings}
\renewcommand*{\chapterpagestyle}{plain.scrheadings}

% abstract
% Same style as table of contents
\section*{Abstract}

\gls{ai} planning generates action sequences that solve complex real-world problems, from robotics to cloud orchestration. Domain-independent planners rely on domain models---typically specified in the \gls{pddl}---to declaratively describe the state space, actions, and goal conditions of a planning task, yet the internal search algorithms of different planners process these models in fundamentally different ways. Prior work has established that the purely textual arrangement of a domain model, without any change in its semantics, can drastically alter planner efficiency. Despite this, existing configuration approaches apply planner-agnostic heuristics that yield inconsistent results: a reordering that accelerates one planner may degrade another.

We investigate whether \glspl{llm} can serve as architecture-aware domain model configurators that restructure PDDL domain descriptions to improve the efficiency of a specific target planner. We propose a multi-stage experimental pipeline comprising four configuration stages: a baseline of original domains, a general-purpose LLM prompt, an architecture-aware prompt that encodes planner-specific structural preferences derived from component-level analysis of each planner's internal architecture, and an LLM-Modulo feedback loop that iteratively refines configurations using quantitative execution telemetry from the target planner. A four-level validation pipeline ensures semantic equivalence throughout.

We evaluate the approach across four architecturally diverse planners (\gls{bfws}, \gls{lama}, DecStar, Madagascar), five \gls{ipc} benchmark domains, and four state-of-the-art LLMs. The full pipeline achieves a statistically significant improvement in planning performance, with architecture-aware prompts nearly doubling the gain of generic prompts. The feedback loop raises the overall improvement rate from 44.1\% to 80.9\% of all configurations and recovers every failure produced by the preceding stage. The magnitude of improvement varies substantially across the four planners, with \gls{lama} and \gls{bfws} exhibiting the largest gains while DecStar and Madagascar prove largely resistant, confirming that responsiveness to domain restructuring is governed by planner architecture. A key finding is that architecture-aware prompting renders the choice of LLM statistically non-significant, demonstrating that the prompt engineering methodology matters more than the underlying model.

These results provide the first empirical validation of planner-aware, iterative LLM-based domain configuration, advancing the field beyond single-pass, planner-agnostic approaches. The findings yield actionable guidelines for mapping structural configurations to target architectures and open a path toward automated, architecture-driven domain engineering in knowledge-intensive AI planning workflows.

\microtypesetup{protrusion=false}

% In case you have trouble with headings reaching into the page numbers, enable the following three lines.
% Hint by http://golatex.de/inhaltsverzeichnis-schreibt-ueber-rand-t3106.html
%
%\makeatletter
%\renewcommand{\@pnumwidth}{2em}
%\makeatother
%
% In case of a strange break in the table of contents,
% a page break can be inserted by issuing the following command at the "right" place in the main text:
%  \addtocontents{toc}{\protect\newpage}
\tableofcontents

\listoffigures

\listoftables

% We use lstlisting environments with caption paramters.
% Thus, we need that command.
% Alternative: \listof{Listing}{List of Listings}
\lstlistoflistings

% mittels \newfloat wurde die Algorithmus-Gleitumgebung definiert.
% Mit folgendem Befehl werden alle floats dieses Typs ausgegeben
%\listof{Algorithmus}{List of Algorithms}
%\listofalgorithms %Ist nur für Algorithmen, die mittels \begin{algorithm} umschlossen werden, nötig

% Abbreviations / Acronyms
\printglossary[type=\acronymtype,title={Abbreviations}]
% \printglossaries
% \printnoidxglossaries
% \printunsrtglossaries cannot be used, because then no indexing happens; source: https://tex.stackexchange.com/a/287128/9075

\microtypesetup{protrusion=true}

% Headline and footline
\renewcommand*{\chapterpagestyle}{scrplain}
\pagestyle{scrheadings}

\chapter{Introduction}\label{sec:introduction}

\section{Context and Motivation}
\label{sec:introduction:context}

% Add Context and Motivation here

\section{Problem Statement and Main Research Question}
\label{sec:introduction:problem-statement}

% Add Problem Statement and Main Research Question here

\subsection{Sub-Questions}
\label{sec:introduction:sub-questions}

% Add Sub-Questions here

\chapter{Background}\label{ch:background}
% Theories and concepts that may be unfamiliar to the audience and are essential for understanding the remainder of the paper

\section{Automated Planning}
\label{sec:bg:planning}

\subsection{Planning Domain Definition Language (PDDL)}
\label{sec:bg:pddl}

\subsection{Domain Models and Configuration Spaces}
\label{sec:bg:domain_models}

\section{AI Planner Architectures}
\label{sec:bg:architectures}

\subsection{Heuristic Forward Search (LAMA / Fast Downward)}
\label{sec:bg:lama}

\subsection{Width-Based Search (BFWS)}
\label{sec:bg:bfws}

\subsection{Decoupled Search (DecStar)}
\label{sec:bg:decstar}

\subsection{SAT-Based Planning (Madagascar)}
\label{sec:bg:madagascar}

\section{Domain Model Configuration}
\label{sec:bg:configuration}

\subsection{Structural Sensitivity of Planners}
\label{sec:bg:sensitivity}

\subsection{Configuration Approaches}
\label{sec:bg:config_approaches}

\section{Large Language Models}
\label{sec:bg:llms}

\subsection{Architecture and Capabilities}
\label{sec:bg:llm_arch}

\subsection{Prompting Strategies}
\label{sec:bg:prompting}

\subsection{LLMs for Planning and Code Transformation}
\label{sec:bg:llms_planning}

\section{The LLM-Modulo Framework}
\label{sec:bg:llm_modulo}

\chapter{Related Work}\label{ch:related_work}
% Identify and describe the key research areas related to your topic

\section{Domain Model Configuration in Classical Planning}
\label{sec:rw:classical_config}

\section{LLMs as Planners}
\label{sec:rw:llms_as_planners}

\section{LLMs as Domain Modellers}
\label{sec:rw:llms_domain_modellers}

\section{LLM-Based Domain Model Configuration}
\label{sec:rw:llm_config}
% Discuss Elis & Georgievski et al.

\section{Research Gap and Positioning}
\label{sec:rw:gap}
% Position your own research within this context; show how it addresses identified gaps

\chapter{Design of the Experimental Pipeline}\label{ch:design}
% Represents the abstract design of the solution

\section{System Overview and Pipeline Architecture}
\label{sec:design:overview}

\section{Stage 0: Baseline Establishment}
\label{sec:design:stage0}

\section{Stage 1: General Prompt Configuration}
\label{sec:design:stage1}

\section{Stage 2: Architecture-Aware Prompt Configuration}
\label{sec:design:stage2}
% Will discuss Prompt Template Design, Planner-Specific Heuristic Rules, and Cross-Testing Protocol here.

\section{Stage 3: LLM-Modulo Feedback Loop}
\label{sec:design:stage3}
% Will discuss Meta-Controller, Diagnostic Feedback, Iterative Refinement, and Termination Criteria here.

\section{Validation Pipeline}
\label{sec:design:validation}

\subsection{V1: PDDL Extraction}
\label{sec:design:v1_extraction}

\subsection{V2: VAL Syntactic Check}
\label{sec:design:v2_val}

\subsection{V3: Identity Check}
\label{sec:design:v3_identity}

\subsection{V4: Semantic Equivalence}
\label{sec:design:v4_semantic}

\section{Computational Safeguards}
\label{sec:design:safeguards}

\chapter{Implementation}\label{ch:implementation}
% Provides implementation details — How is the design practically realised?

\section{System Architecture and Codebase}
\label{sec:impl:codebase}
% Python codebase structure, module organisation
% How the pipeline stages are orchestrated programmatically

\section{Docker Containerisation}
\label{sec:impl:docker}
% How the C++ planners are wrapped in Docker containers
% Ensures consistent, reproducible execution across platforms
% Isolates fragile planner dependencies from the host system

\section{Planner Integration}
\label{sec:impl:planners}
% How each planner is integrated into the pipeline
% Described in the context of integration (command-line interfaces, output parsing, Docker images)
% NOT abstract background — the reader encounters planners here after understanding the pipeline design

\subsection{LAMA (Fast Downward)}
\label{sec:impl:lama}
% Heuristic forward search planner
% Integration specifics: Fast Downward build, command-line flags, output format

\subsection{BFWS}
\label{sec:impl:bfws}
% Width-based search planner
% Integration specifics: build process, novelty parameters, output parsing

\subsection{DecStar}
\label{sec:impl:decstar}
% Decoupled/factored search planner
% Integration specifics: factoring configuration, Docker setup

\subsection{Madagascar}
\label{sec:impl:madagascar}
% SAT-based planner
% Integration specifics: SAT encoding parameters, bounded horizon, output format

\section{LLM Integration}
\label{sec:impl:llm_integration}
% LiteLLM API abstraction layer
% How models are called uniformly across providers (OpenAI, Anthropic, Google, Meta)
% Prompt injection, response parsing, error handling

\chapter{Experimental Setup}\label{ch:experimental_setup}
% Defines the specific ingredients, rules, and measurement instruments of the experiment

\section{Selection of Domains}
\label{sec:setup:domains}
% The 5 IPC benchmark domains: Depots, Snake, Visitall, Barman, Ricochet-Robots
% Selection criteria: mix of old/new, different difficulty levels, structural diversity

\section{Selection of Planners}
\label{sec:setup:planners}
% The 4 planners and their exact versions
% Selection methodology and rationale (covering 4 distinct search paradigms)

\section{Selection of LLMs}
\label{sec:setup:llms}
% The 4 LLMs: GPT-4o, Claude 3.5 Sonnet, Gemini 1.5 Pro, Llama 3
% Selection criteria: mix of proprietary/open-source, reasoning/coding specialisation

\section{Hardware and Environment}
\label{sec:setup:hardware}
% Physical specs: Intel i7-13650HX, 16GB RAM, OS
% Local dedicated machine to prevent noisy-neighbour variance

\section{Hyperparameters and Reproducibility}
\label{sec:setup:hyperparameters}
% LLM temperature = 0.0, fixed random seeds
% Planner timeout/memout limits
% Number of runs per configuration
% Measures taken to ensure deterministic reproducibility

\section{Evaluation Metrics}
\label{sec:setup:metrics}
% IPC Score: official International Planning Competition scoring formula
% Coverage: how many problem instances were solved
% PAR10: penalised average runtime for unsolved instances
% Runtime and Plan Cost: speed and efficiency of the produced plan
% Validation Rates: percentage of LLM outputs that pass the 4-level validation pipeline

\section{Statistical Methods}
\label{sec:setup:stats}
% Friedman Test: proves at least one LLM/stage is significantly different
% Wilcoxon Signed-Rank Test: pairwise comparison (Stage 2 vs. Stage 1)
% Kruskal-Wallis Test: comparison across multiple groups
% Cliff's Delta: effect size — not just "if" better, but "how much" better
% Bonferroni / Nemenyi post-hoc corrections for multiple comparisons

\chapter{Results}\label{ch:results}
% Pure results presentation and interpretation — no design or methodology

\section{Baseline Characterisation (Stage 0)}
\label{sec:results:stage0}
% Tables showing planner performance on original, unmodified domains
% Interpretation: which domains are inherently hard, which planners struggle the most
% Establishes the reference point for all subsequent comparisons

\section{General Prompt Results (Stage 1)}
\label{sec:results:stage1}
% Did a simple "improve this code" prompt work?
% Validation rates, IPC score gains, coverage changes
% Analysis: worked partially, but often generated broken or unhelpful code

\section{Architecture-Aware Prompt Results (Stage 2)}
\label{sec:results:stage2}
% The breakthrough results
% Validation rates, IPC score jumps, coverage improvements
% Cross-testing analysis: proving architectural rules actually matter
% Comparison with Stage 1: interpreting why architecture-aware hints outperform generic ones

\section{Feedback Loop Results (Stage 3)}
\label{sec:results:stage3}
% Validation and convergence analysis
% How the loop caught errors and fixed them iteratively
% Success rate improvement (e.g., ~44\% to ~80\%)
% Iterative learning patterns and failure recovery analysis

\section{Cross-Stage Comparative Analysis}
\label{sec:results:cross_stage}
% Full pipeline progression: Stages 0 → 1 → 2 → 3 side-by-side
% Ablation study isolating the value of each stage
% LLM effectiveness comparison: which model was smartest
% Planner responsiveness: which planner benefited most from good code
% Domain sensitivity: which domains responded best
% Token efficiency and cost analysis

\section{Statistical Validation}
\label{sec:results:statistical_validation}
% Omnibus tests (Friedman): proving overall significance
% Pairwise comparisons (Wilcoxon): Stage 2 vs. Stage 1, Stage 3 vs. Stage 2
% Effect sizes (Cliff's Delta): magnitude of improvements
% Proving to the committee that the results are not coincidence

\section{Comparison with Prior Work}
\label{sec:results:prior_work}
% Direct comparison with Daniel Elis's thesis results
% How the architecture-aware, feedback-loop pipeline outperforms the generic baseline

\chapter{Threats to Validity}\label{ch:threats}
% Dedicated chapter discussing potential threats following the standard four-type framework

\section{Internal Validity}
\label{sec:threats:internal}
% Factors that could compromise causal conclusions within the experiment
% Hardware noise: mitigated by dedicated local machine, but single-machine variance remains
% API non-determinism: LLM outputs may vary despite temperature=0.0 due to provider-side batching
% Docker overhead: containerisation adds latency that may affect absolute runtime measurements
% Order effects: whether the sequence of problem instances affects planner warm-up behaviour

\section{External Validity}
\label{sec:threats:external}
% Generalisability of findings beyond this specific experiment
% Limited to 5 IPC benchmark domains — results may not generalise to all PDDL structures
% Limited to 4 planners — other planner architectures (e.g., temporal, HTN) are untested
% STRIPS subset of PDDL — does not cover numeric fluents, temporal planning, or ADL features
% LLM model versions: results are tied to specific model snapshots that may be deprecated

\section{Construct Validity}
\label{sec:threats:construct}
% Whether the measurements actually capture the intended constructs
% IPC score as a proxy for "planner efficiency" — may not capture all dimensions of planning quality
% "Architecture-awareness" operationalised as manually crafted rules — other operationalisations are possible
% Semantic equivalence check (V4) relies on structural comparison, not formal verification
% Prompt design decisions may introduce experimenter bias

\section{Statistical Conclusion Validity}
\label{sec:threats:statistical}
% Reliability of statistical inferences drawn from the data
% Sample size per experimental cell: whether sufficient for non-parametric test power
% Multiple comparisons: risk of Type I errors across the many pairwise tests conducted
% Non-parametric test assumptions: ordinal data requirements, independence assumptions
% Effect size interpretation: Cliff's Delta thresholds are conventional, not domain-specific

\chapter{Conclusion and Outlook}\label{ch:conclusion}

\section{Summary of Contributions}
\label{sec:concl:summary}
% State the most important outcomes:
% 1. Introduced a 4-stage pipeline for architecture-aware domain model configuration
% 2. Proved that architecture-aware prompting mathematically outperforms generic prompting
% 3. Demonstrated that a closed feedback loop increases validation/improvement rates significantly
% 4. Showed that the engineering methodology matters more than the choice of LLM

\section{Answers to Research Questions}
\label{sec:concl:rq_answers}
% Copy the research questions from Chapter 1 and provide definitive, data-backed answers
% Each question gets a clear response grounded in the findings from Chapter 7

\section{Limitations}
\label{sec:concl:limitations}
% Software and practical limitations ONLY (experimental validity is in Chapter 8)
% API latency: feedback loop takes several minutes per domain
% Cost: token consumption across multiple LLM calls
% Manual rule engineering: architecture-aware rules were hand-crafted from planner source code
% Single-machine execution: limits the scale of experiments that can be run

\section{Future Work}
\label{sec:concl:future_work}
% Each limitation should spawn an idea for future work
% Automated rule extraction: use a separate LLM agent to read planner C++ source code and extract rules
% Scaling to more planners/domains: temporal, HTN, numeric PDDL
% Online adaptation: real-time feedback during planning, not just offline batch processing
% Fine-tuning LLMs specifically for PDDL configuration tasks
% Integration into planning competition infrastructure


%%% ===============================================================================
%%% Bibliography
%%% ===============================================================================

\printbibliography

% Enfore empty line after bibliography
\ \\
%

\clearpage
\appendix
\addcontentsline{toc}{part}{Appendix}

\chapter{Architecture-Aware Prompts}\label{ch:appendix:prompts}
%%% ===============================================================================

This appendix provides the complete, verbatim text of the four architecture-aware prompts used in Stage~2 of the experimental pipeline, as discussed in \Cref{sec:setup:arch_aware_prompt}. Each prompt is tailored to the internal mechanics of its target planner.

\section{LAMA-First Architecture-Aware Prompt}

\begin{lstlisting}[language={}, caption={Architecture-aware prompt for LAMA-first.}, label={lst:prompt:lama}, basicstyle=\small\ttfamily, numbers=none, frame=single, breaklines=true]
You are a PDDL domain engineering expert specializing in optimizing domain files
for the LAMA-first planning engine (Fast Downward framework). You possess deep
knowledge of how LAMA's internal architecture - specifically its Translation module
(grounding + invariant synthesis), Knowledge Compilation module (DTG construction),
and Search module (multi-heuristic FF/add + Landmark heuristic with 4-queue
preferred operator system) - processes PDDL input.

Your task is to reorder the following PDDL domain file to improve LAMA's planning
efficiency (coverage, runtime, IPC score).

STRICT CONSTRAINTS:
- Do NOT add, remove, or rename any predicates, types, constants, or actions.
- Do NOT change the logical meaning of any precondition or effect.
- Do NOT alter parameter types or add/remove parameters.
- Do NOT add comments, explanations, or formatting changes.
- Return ONLY the complete, valid, reordered PDDL domain file.

ARCHITECTURE CONTEXT:
LAMA processes PDDL through a three-stage pipeline:
1. TRANSLATION: Grounds all action schemas, discovers mutex groups (mutually
   exclusive fact groups) via invariant synthesis, and constructs a Finite-Domain
   Representation (SAS+/FDR). Mutex information is critical because it directly
   determines the quality of landmark ordering constraints used later in search.
2. KNOWLEDGE COMPILATION: Builds Domain Transition Graphs (DTGs) for each state
   variable, constructs fast operator-applicability lookup structures, and prepares
   axiom evaluation mechanisms.
3. SEARCH: Runs greedy best-first search (for initial plan), then iterated weighted
   A* with restarts. Uses TWO heuristics simultaneously in separate queues:
   - FF/add heuristic (delete-relaxation based, additive cost propagation)
   - Landmark heuristic (counts unachieved mandatory waypoints)
   Each heuristic generates PREFERRED OPERATORS - actions deemed especially promising - 
   which receive priority via a +1000 boost mechanism in a 4-queue system.

Because of this architecture, the textual ordering of predicates affects mutex
discovery and landmark ordering quality. The ordering of actions affects successor
generation, tie-breaking, and which operators are evaluated first. The ordering of
preconditions directly affects which "best support" is selected during FF's
relaxed-plan extraction, determining WHICH actions become preferred operators.
The ordering of effects affects delete-relaxed forward propagation and landmark
acceptance tracking.

ARCHITECTURE-AWARE REORDERING RULES (apply these in order of priority):

1. PREDICATE DECLARATIONS: Place goal-relevant predicates (those describing
   positional state, object relationships, or task completion) at the top of the
   :predicates block. Follow with frequently-occurring predicates. Place static
   predicates (those never modified by any :effect) last.
   REASON: Early placement accelerates mutex discovery during invariant synthesis,
   improving landmark ordering constraint quality.

2. ACTION ORDERING: Place actions with fewer preconditions earlier in the domain
   file. Among actions with similar precondition counts, prioritize "enabling"
   actions - those whose effects create preconditions for other actions.
   REASON: LAMA's greedy best-first search evaluates operators in listing order
   during deferred evaluation. Simpler, enabling actions listed first produce
   better initial plans during the speed-first greedy sprint.

3. PRECONDITION ORDERING: Within each action's :precondition block:
   a) Place STATIC predicates (never changed by any effect) first.
   b) Then place the MOST RESTRICTIVE conditions (those satisfied in fewest states).
   c) Then group remaining predicates by relatedness (same predicate name or
      referencing the same object types).
   REASON: Static predicates first enable fast applicability short-circuiting.
   Restrictive conditions first cause earlier failure detection. The ordering
   also affects FF's relaxed-plan backward trace, which selects "best supports"
   and determines preferred operators.

4. EFFECT ORDERING: Within each action's :effect block:
   a) Place ADD-EFFECTS first (positive effects without "not").
   b) Within add-effects, place goal-relevant effects before auxiliary effects.
   c) Place DELETE-EFFECTS (negative effects with "not") last.
   REASON: LAMA's FF/add heuristic operates in a delete-relaxed world where only
   add-effects drive cost propagation. Goal-relevant add-effects first accelerate
   the forward propagation phase and improve landmark acceptance tracking.

5. PARAMETER CONSISTENCY: Maintain consistent parameter ordering across all action
   schemas (e.g., always agent before object before location).
   REASON: Consistent ordering simplifies FDR translation and supports causal graph
   structural consistency.

DOMAIN TO REORDER:
{domain_str}
\end{lstlisting}

\section{DecStar Architecture-Aware Prompt}

\begin{lstlisting}[language={}, caption={Architecture-aware prompt for DecStar.}, label={lst:prompt:decstar}, basicstyle=\small\ttfamily, numbers=none, frame=single, breaklines=true]
You are a PDDL domain engineering expert specializing in optimizing domain files
for the DecStar planning engine. You possess deep knowledge of how DecStar's 
internal architecture - specifically its h^2 preprocessing, Star Topology factoring
engine (Center/Leaf assignment via Causal Graph and InteractionGraph analysis),
and Decoupled Search (Center Actions vs. silent Leaf Actions with Pricing
Functions) - processes PDDL input.

Your task is to reorder the following PDDL domain file to improve DecStar's 
planning efficiency (coverage, runtime, IPC score).

STRICT CONSTRAINTS:
- Do NOT add, remove, or rename any predicates, types, constants, or actions.
- Do NOT change the logical meaning of any precondition or effect.
- Do NOT alter parameter types or add/remove parameters.
- Do NOT add comments, explanations, or formatting changes.
- Return ONLY the complete, valid, reordered PDDL domain file.

ARCHITECTURE CONTEXT:
DecStar extends Fast Downward with DECOUPLED STATE-SPACE SEARCH. After standard
FD translation (grounding + FDR) and aggressive h^2 mutex pruning, DecStar builds
a STAR TOPOLOGY that separates state variables into:
- A CENTRAL HUB ("Center"): highly-connected, entangled variables
- Independent LEAF modules: isolated subsystems that can only interact through
  the Center

The factoring engine enforces a GOLDEN RULE: Leaves cannot interact with other
Leaves directly. It uses the Causal Graph to identify Strongly Connected
Components (SCCs), strips off isolated "tips" as Leaves, and whatever tangled
mass remains becomes the Center.

CRITICAL PERFORMANCE MECHANISM:
- CENTER ACTIONS (affecting any Center variable): Branch the main search tree,
  generating new states in the open list  -  EXPENSIVE.
- LEAF ACTIONS (affecting only one Leaf): Handled silently by localized Dijkstra
  within the current state  -  CHEAP.

The MORE actions that qualify as Leaf-only operators, the greater the combinatorial
compression. Actions with preconditions/effects spanning multiple subsystems force
Center classification and defeat decoupling.

SAFETY CONSTRAINTS: Factoring has a 30-second timer and a maximum leaf size limit.
If the topology cannot be built, DecStar falls back to standard explicit-state
search, losing all decoupling benefits.

ARCHITECTURE-AWARE REORDERING RULES (apply these in order of priority):

1. PREDICATE DECLARATIONS: Group predicates by the object types they describe.
   Place predicates involving SINGLE object types first (these become candidate
   leaf variables). Place predicates describing INTERACTIONS between different
   object types last (these likely become center variables).
   REASON: Cleaner predicate grouping leads to more efficient mutex group
   construction, which feeds directly into the SCC analysis that determines
   the Center/Leaf partition.

2. ACTION ORDERING: Place actions that operate on a SINGLE independent subsystem
   (e.g., manipulating one package, one block, one room) FIRST. Place actions
   that COORDINATE across multiple subsystems (e.g., loading a package onto a
   truck, connecting two rooms) LAST.
   REASON: This ordering emphasizes the hub-and-spoke structure that DecStar's 
   Fork/Inverted-Fork factoring algorithms are designed to detect, helping the
   factoring engine identify valid topologies within the 30-second timer.

3. PRECONDITION ORDERING: Within each action, place preconditions referencing
   the action's PRIMARY manipulated object first. Place preconditions referencing 
   SHARED or COORDINATING state (hub variables, transport vehicles, etc.) last.
   REASON: The InteractionGraph analyzes precondition-effect dependencies to
   detect cross-leaf violations. Local-first ordering makes this analysis more
   efficient.

4. EFFECT ORDERING: Place effects modifying the action's PRIMARY (local) object 
   first. Place effects modifying shared or coordinating state last.
   REASON: Leaf-local effects first enables faster pricing function computation
   during silent Dijkstra execution within leaf notebooks.

5. GLOBAL PREDICATE DEMOTION: If any predicates appear in preconditions or effects
   of nearly ALL actions (e.g., total-cost counters, phase flags), place them at
   the END of the :predicates block and at the END within each precondition/effect
   block where they appear.
   REASON: Global predicates draw causal edges connecting all objects, potentially
   collapsing the entire task into one SCC and destroying the Star Topology.
   Demoting their textual position may influence variable ordering in the
   translator, mitigating this destructive effect.

6. PARAMETER CONSISTENCY: Maintain consistent parameter ordering, with the
   primary manipulated object as the first parameter in all actions.
   REASON: Consistent ordering supports cleaner causal graph construction.

DOMAIN TO REORDER:
{domain_str}
\end{lstlisting}

\section{BFWS-Preference Architecture-Aware Prompt}

\begin{lstlisting}[language={}, caption={Architecture-aware prompt for \gls{bfws}-Preference.}, label={lst:prompt:bfws}, basicstyle=\small\ttfamily, numbers=none, frame=single, breaklines=true]
You are a PDDL domain engineering expert specializing in optimizing domain files
for the LAPKT-BFWS-Preference planning engine (IPC 2018 Agile Track winner). You
possess deep knowledge of how BFWS-Preference's internal architecture - specifically
its DIRECT MEMORY INJECTION parsing (no intermediate file), BUCKET SORTING via
partition = (1000 x #g) + #r, NOVELTY evaluation w(s), and its TWO-ENGINE RELAY
(LAPKT then Fast Downward) - processes PDDL input.

Your task is to reorder the following PDDL domain file to improve BFWS-Preference's
planning efficiency (coverage, runtime, IPC score).

STRICT CONSTRAINTS:
- Do NOT add, remove, or rename any predicates, types, constants, or actions.
- Do NOT change the logical meaning of any precondition or effect.
- Do NOT alter parameter types or add/remove parameters.
- Do NOT add comments, explanations, or formatting changes.
- Return ONLY the complete, valid, reordered PDDL domain file.

ARCHITECTURE CONTEXT:
BFWS-Preference is a TWO-ENGINE RELAY system:

ENGINE 1 (LAPKT - The Sprinter):
The Python manager DIRECTLY INJECTS parsed PDDL into C++ memory via bindings
(add_atom, add_precondition, add_effect). There is NO intermediate .sas file.
C++ internal ID numbers are assigned in the EXACT ORDER elements appear in the
PDDL text. These IDs control tie-breaking EVERYWHERE in the engine.

The search uses BUCKET SORTING: partition = (1000 x #g) + #r, where:
- #g = number of MISSING GOALS (simple count)
- #r = number of RELEVANT FLUENTS from a delete-free relaxed plan (H_Add_Rp_Fwd)

CRITICAL: H_Add_Rp_Fwd breaks ties among equally-optimal relaxed actions by
selecting the one with the LOWEST INTERNAL ID NUMBER. This means changing the
textual order of actions in the PDDL file DIRECTLY CHANGES which relaxed actions
are selected, which changes the #r count, which changes the BUCKET ASSIGNMENT,
which changes the NOVELTY score w(s), which completely reroutes the search tree.

ENGINE 2 (Fast Downward/LAMA - The Marathon Runner):
After Engine 1 finds an initial plan, Engine 2 re-parses the SAME PDDL file using
standard Fast Downward translation (with h^2 preprocessing), then runs weighted A*
with LAMA heuristics to find cheaper plans within the time limit.

ARCHITECTURE-AWARE REORDERING RULES (apply in order of priority):

1. ACTION ORDERING [CRITICAL  -  HIGHEST IMPACT]:
   Place actions whose effects directly ACHIEVE GOAL predicates FIRST in the
   domain file. Then place actions that CREATE PRECONDITIONS for goal-achieving
   actions. Place general utility actions last.
   REASON: Goal-achieving actions receive the lowest C++ IDs. When H_Add_Rp_Fwd
   breaks ties during relaxed plan extraction, it selects the lowest-ID action,
   which will now be goal-relevant. This produces tighter #r values and more
   discriminating bucket assignments, improving novelty-based search focus.

2. PREDICATE DECLARATIONS: Place GOAL-RELEVANT predicates (describing task
   completion, object positioning, desired states) first. Then place action-
   enabling predicates. Place auxiliary/static predicates last.
   REASON: Predicate IDs are also assigned sequentially. Goal-relevant predicates
   with lower IDs are prioritized in the novelty tuple tables, improving novelty
   discrimination.

3. PRECONDITION ORDERING: Within each action, place the most GOAL-PROXIMAL
   predicates first (those appearing in goal conditions or as effects of goal-
   achieving actions). Place static/auxiliary predicates last.
   REASON: Precondition order affects the H_Add_Rp_Fwd backward trace,
   influencing which support actions are selected for the relaxed plan. Goal-
   proximal preconditions first guide the trace toward goal-relevant fluents.

4. EFFECT ORDERING: Place effects that ACHIEVE or ENABLE goal predicates first.
   Place "cleanup" or state-removal effects last.
   REASON: During delete-free forward reachability, goal-relevant add-effects
   first ensure early discovery of goal-adjacent facts, reducing #r and producing
   tighter partitions.

5. DUAL-ENGINE COMPATIBILITY: Ensure the reordering is also compatible with Fast
   Downward's pipeline (Engine 2). Static predicates should be identifiable early
   in precondition blocks for FD's invariant synthesis. Add-effects before
   delete-effects for FD's relaxed planning.

DOMAIN TO REORDER:
{domain_str}
\end{lstlisting}

\section{Madagascar Architecture-Aware Prompt}

\begin{lstlisting}[language={}, caption={Architecture-aware prompt for Madagascar (MpC).}, label={lst:prompt:madagascar}, basicstyle=\small\ttfamily, numbers=none, frame=single, breaklines=true]
You are a PDDL domain engineering expert specializing in optimizing domain files
for the Madagascar (MpC configuration) planning engine. You possess deep knowledge
of how Madagascar's internal architecture - specifically its CNF encoding of PDDL
into Boolean SAT formulas using E-step parallel plans, its fixpoint binary mutex
(invariant) discovery, its backward-chaining variable selection heuristic, and its
CDCL (Conflict-Driven Clause Learning) solver with Unit Propagation - processes
PDDL input.

Your task is to reorder the following PDDL domain file to improve Madagascar's
planning efficiency (coverage, runtime, IPC score).

STRICT CONSTRAINTS:
- Do NOT add, remove, or rename any predicates, types, constants, or actions.
- Do NOT change the logical meaning of any precondition or effect.
- Do NOT alter parameter types or add/remove parameters.
- Do NOT add comments, explanations, or formatting changes.
- Return ONLY the complete, valid, reordered PDDL domain file.

ARCHITECTURE CONTEXT:
Madagascar is a SAT-BASED planner. It does NOT search through states like
traditional planners. Instead, it:
1. ENCODES the entire PDDL problem into a massive Boolean formula (Phi_t) in
   Conjunctive Normal Form (CNF). Each PDDL predicate at each time step becomes
   a True/False variable. Each PDDL action at each time step becomes a True/False
   variable. The formula represents ALL constraints simultaneously.
2. Discovers BINARY MUTEXES (invariants) via fixpoint computation  -  rules like
   "package cannot be at LocationA AND LocationB simultaneously." These are stored
   ONCE in memory and dynamically linked across time steps, providing massive
   memory savings.
3. Uses E-step (existential-step) PARALLEL PLANS that allow multiple non-
   interfering actions in the same time step, dramatically shrinking formula size.
4. Runs CONCURRENT SAT SOLVERS at multiple time horizons (Algorithm C for MpC:
   exponential steps with equal CPU allocation).
5. Each SAT solver uses CDCL with a BACKWARD-CHAINING HEURISTIC: instead of
   generic conflict-driven guessing, it starts from GOAL variables, identifies
   which action variables directly achieve those goals, and prioritizes evaluating
   them first. Sub-goals then dynamically become the next priority.

CRITICAL: The order of PDDL elements determines the internal SAT variable IDs.
The backward-chaining heuristic's priority queue processes variables by their IDs.
Goal-achieving actions and goal predicates with LOWER IDs are evaluated FIRST.

ARCHITECTURE-AWARE REORDERING RULES (apply in order of priority):

1. PREDICATE DECLARATIONS: Place GOAL predicates (those describing desired end
   states) FIRST. Then place predicates that are effects of goal-achieving actions.
   Group mutually exclusive predicates together (e.g., all positional predicates
   for the same object type adjacently, like at-locA, at-locB, at-locC).
   REASON: Goal predicates receive the lowest SAT variable IDs, placing them first
   in the backward-chaining priority queue. Adjacent mutually-exclusive predicates
   are discovered as binary mutex pairs sooner during fixpoint iteration.

2. ACTION ORDERING: Place actions that DIRECTLY ACHIEVE goal predicates FIRST.
   Then place actions that achieve subgoals (preconditions of goal-achieving
   actions). Place auxiliary/setup actions last.
   REASON: The backward-chaining heuristic starts from goal variables and
   identifies their causal actions. Lower SAT variable IDs for goal-achieving
   actions place them higher in the priority queue, focusing the CDCL solver's
   initial guesses on the most productive variables.

3. PRECONDITION ORDERING: Place EASILY SATISFIABLE preconditions first (those
   true in many states or established by simple actions). Place preconditions
   requiring complex multi-step setup last.
   REASON: Preconditions are encoded as CNF implications. Easily-satisfiable
   conditions trigger longer UNIT PROPAGATION domino chains, forcing more
   variables automatically and reducing the number of explicit guesses needed.

4. EFFECT ORDERING: Place ADD-EFFECTS first, then DELETE-EFFECTS. Within
   add-effects, place GOAL-RELEVANT effects first.
   REASON: Although SAT treats literals symmetrically, add-effects placed first
   align with the backward-chaining heuristic's goal-focused search direction
   and produce frame axiom clauses in an order that supports faster propagation.

5. MUTUAL EXCLUSION ALIGNMENT: In the :predicates block, place predicates
   describing the SAME PROPERTY of the same object type ADJACENTLY (e.g., group
   all location predicates for trucks together, all status predicates for
   packages together).
   REASON: Madagascar's fixpoint invariant finder discovers binary mutexes by
   checking predicate pairs. Adjacent same-property predicates are checked
   earlier, allowing critical invariants to be discovered sooner.

6. PARAMETER CONSISTENCY: Use consistent parameter ordering across all actions.
   REASON: Regular grounding patterns produce more uniform CNF clause structures
   that the CDCL solver exploits through learned clause pattern recognition.

DOMAIN TO REORDER:
{domain_str}
\end{lstlisting}

\pagestyle{empty}
\renewcommand*{\chapterpagestyle}{empty}
\Versicherung
\end{document}
